% Standalone research fragment. Dependencies: the general terminal-path gap,
% and the separately audited saturated-ring bound.
\section{Rings with arbitrary four-terminal blocks}
\label{sec:general-four-pole-rings}

Let $K_1,\ldots,K_k$ be disjoint connected loopless multigraphs.
Each block has two distinct left ports and two distinct right ports.
The four ports have degree two, and all other vertices have degree three.
Join consecutive blocks by a matching of size two between their corresponding
right and left ports. Indices are cyclic. Put
\[
 h_i=\beta(K_i),\qquad H=\sum_{i=1}^k h_i,\qquad
 W=H+k,\qquad r=\beta(G)=W+1.
\]
The degree sum gives $|V(K_i)|=2h_i+2$, so $h_i\ge1$.
A block is \emph{cross-saturated} if every left-to-right terminal pair has
exactly $2^{h_i}$ simple paths.

\begin{proposition}[Arbitrary rings with at least eleven blocks]
\label{prop:general-four-pole-ring}
Suppose $k\ge11$. If at least one block is not cross-saturated, then
\[
 c(G)<2^{r-1}.
\]
Consequently every such ring satisfies
\[
 c(G)\le 2^{r-1}+20r2^{r/2}.
\]
The second conclusion uses the saturated-ring theorem when every block is
cross-saturated.
\end{proposition}

\begin{proof}
Let $P_i$ be the $2\times2$ matrix of cross-path counts in $K_i$.
Every entry is at most $2^{h_i}$.
If block $j$ is not cross-saturated, the terminal-path gap gives
\[
 \sum_{a,b}(P_j)_{ab}\le \left(4-\frac1{16}\right)2^{h_j}.
\]
Fixed interface matchings can be incorporated by permutation matrices.
Cycles using one edge of every interface are counted by the corresponding
matrix-product trace.
Replacing every factor except $P_j$ by its entrywise upper bound gives
\[
 c_{\mathrm{odd}}(G)
 \le 2^{H-h_j}2^{k-2}\sum_{a,b}(P_j)_{ab}
 \le \frac{63}{64}2^W.
\]

All interface intersection sizes have the same parity.
Indeed, each block has an even number of cycle edges in its boundary.
It remains to bound cycles using zero or two edges of every interface.
Cycles internal to one block contribute at most
\[
 \sum_i 2^{h_i}\le k2^{H-k+1}.
\]

If the used interfaces form a proper nonempty interval, the cycle meets a
contiguous interval of $s$ blocks, where $2\le s\le k$.
The case $s=k$ omits one interface and uses all blocks.
It uses a cap in each end block and a through pair in every intermediate block.
Each cap count is at most $2^{h_i}$ by the affine terminal-path bound.
Each through-pair count is at most $2^{h_i}$ by the affine four-terminal
join count. For each $s$, there are $k$ possible block intervals.
The omitted blocks have total rank at least $k-s$.
Thus their contribution is at most the first expression below:
\[
 k\sum_{s=2}^{k}2^{H-k+s}<2k2^H.
\]

For the remaining case, every interface is used twice.
Write $D_i$ for the number of through pairs in block $i$.
Write $C_i$ for its number of disjoint left and right caps.
These are precisely the three pairing types of four-terminal forests.
Their total satisfies
\[
 D_i+C_i\le2^{h_i}.
\]
A single global cycle uses through pairs everywhere, or double caps in
exactly one block. Two double-cap blocks would separate the selected edges.
Therefore this case contributes at most
\[
 \prod_iD_i+\sum_i C_i\prod_{j\ne i}D_j
 \le\prod_i(D_i+C_i)\le2^H.
\]

Combining the even cases gives
\[
 \frac{c_{\mathrm{even}}(G)}{2^W}
 <(2k+1)2^{-k}+k2^{-2k+1}<\frac1{64}
 \qquad(k\ge11).
\]
The final expression decreases with $k$.
At $k=11$, it equals $23/2048+22/4194304<1/64$.
Adding the odd contribution proves the strict bound.
The saturated case follows from the separately proved saturated-ring bound.
\end{proof}

\paragraph{Scope.}
The proposition permits arbitrary block ranks and arbitrary cross-path defects.
It reduces this construction class to saturated rings and rings with at most
ten blocks. It does not bound arbitrary graphs by the half-exponent error.
